\section{The Onsager tensor and its exact degeneracy}
\label{ons:section}

Moment matching separates the macroscopic parameters from the microscopic
deviation. To determine the first-order contribution of this deviation to
the flux, we project the Euler transport forcing off the null space,
solve for its response using the weighted collision form, and take the
five moment fluxes. This defines the Onsager tensor for the cube-truncated model.
Write $\theta=(\alpha,\beta,\gamma)\in\Theta_+=\mathcal O$, with
$\beta\in\R^3$, and
$N=(\alpha+\beta\cdot k+\gamma|k|^2)^{-1}$.
All uniform statements are made on a fixed compact set
$\Theta\Subset\Theta_+$. Reciprocity is a basic feature of Onsager theory
\cite{Onsager1931}; here its definition, null directions and transport
properties follow from the four-leg dissipation form of this model.

\subsection{The weighted cell problem and transport tensor}

Set $\kappa_a=N\Psi_a$, $a=0,\ldots,4$, and retain the projections
$P_N,Q_N$ from Proposition \ref{geom:coercivity}.
These $\kappa_a$ are the natural generators. With the Gram matrix $M$
in \eqref{hyd:matrices}, the projection is
\[
 P_NF=(N\Psi)^TM^{-1}\int_DN\Psi F\,\dd k.
\]
This is the same projection as that obtained from the orthonormal basis
in Section \ref{geom:section}. For $F\in Q_NL^\infty(D)$, write
$L_{N,Q}^{-1}F$ for the normalized solution of
\begin{equation}\label{ons:weighted-cell}
 u\in Q_NH_\nu,\qquad
 \mathfrak a_N(u,h)=\langle F,h\rangle
 \quad(h\in H_\nu).
\end{equation}
The pairing on the right is the unweighted integral pairing.

\begin{lemma}[Weighted inverse and smooth parameter dependence]\label{ons:cell}
 Equation \eqref{ons:weighted-cell} has a unique solution, and
 \begin{equation}\label{ons:cell-bound}
 \|u\|_{H_\nu}+\|\nu_*u\|_\infty\le C_\Theta\|F\|_\infty.
 \end{equation}
 If $F=F_\theta\in Q_NL^\infty$ depends smoothly on $\theta$, then
 $u_\theta$ depends smoothly on $\theta$ with values in
 $\mathcal Y=\{u:\nu_*u\in L^\infty\}$. Each derivative of fixed order
 is uniformly bounded on compact parameter sets.
 For $F,H\in Q_NL^\infty(D)$, the following pairing is well defined:
 \begin{equation}\label{ons:inverse-pairing}
 \langle F,L_{N,Q}^{-1}H\rangle
 =\mathfrak a_N(L_{N,Q}^{-1}F,L_{N,Q}^{-1}H).
 \end{equation}
\end{lemma}
\begin{proof}
 Proposition \ref{geom:cell} gives the solution and
 \eqref{ons:cell-bound}. For every $h\in H_\nu$,
 \[
 |\langle F,h\rangle|
 \le\|F/\sqrt{\nu_*}\|_2\|h\|_{H_\nu}
 \le C\|F\|_\infty\|h\|_{H_\nu},
 \]
 since $\nu_*^{-1}\in L^1$. The equation initially obtained for bounded
 test functions therefore extends continuously to all of $H_\nu$.
 Testing with another weighted solution proves
 \eqref{ons:inverse-pairing}.

 For parameter dependence, use the reference state $n_*$, set
 $Q_*=I-P_{n_*}$, and fix the null space with $U_N$ from
 Lemma \ref{geom:frame}. The operators
 $\widetilde L_N=U_N^*L_NU_N$ form a smooth family of bounded
 isomorphisms between the fixed spaces
 $Q_*\mathcal Y$ and $Q_*\mathcal X$. Their inverse norms are uniformly
 controlled by \eqref{ons:cell-bound}. The identity
 \[
 \partial_\theta\widetilde L_N^{-1}
 =-\widetilde L_N^{-1}(\partial_\theta\widetilde L_N)
       \widetilde L_N^{-1}
 \]
 and its successive derivatives give the claim. Conjugating back by
 $U_N$ returns to the original spaces. The solutions remain in weighted
 spaces throughout; an unweighted spectral gap is not needed.
\end{proof}

Define the fifteen forcing terms, their responses, and the transport tensor by
\begin{equation}\label{ons:definition}
 \begin{aligned}
 \mathcal G_j^a&=Q_N(v_j\kappa_a),\qquad
 u_j^a=L_{N,Q}^{-1}\mathcal G_j^a,\\
 \mathsf K_{ij}^{ab}(N)&=\langle\mathcal G_i^a,u_j^b\rangle,
 \qquad 1\le i,j\le3,\quad0\le a,b\le4.
 \end{aligned}
\end{equation}
The index $j$ specifies a spatial direction, and $a$ specifies mass,
one of the three momentum components, or energy. The fifteen forcing terms
thus correspond to variations of five parameters in three spatial
directions. The projection removes the invariant component, the weighted
inverse gives the microscopic response, and the final pairing extracts
its contribution to flux $(i,a)$. The forcing terms are real and bounded.
The five-dimensional Gram formula for the projection and
Lemma \ref{ons:cell} show that the tensor is smooth in $\theta$.

\begin{proposition}[Reciprocity and nonnegativity]\label{ons:tensor}
 The tensor \eqref{ons:definition} is a real symmetric positive
 semidefinite $15\times15$ matrix:
 \begin{equation}\label{ons:reciprocity}
 \mathsf K_{ij}^{ab}=\mathsf K_{ji}^{ba},\qquad
 \sum_{i,j,a,b}\xi_i^a\mathsf K_{ij}^{ab}\xi_j^b\ge0
 \quad(\xi\in\R^{3\times5}).
 \end{equation}
 Each parameter derivative of fixed order is uniformly bounded on
 compact positive parameter sets. For a real gradient array $\xi$, set
 $F_\xi=\sum_{j,a}\mathcal G_j^a\xi_j^a$. Then
 \begin{equation}\label{ons:variational}
 \xi^{\mathsf T}\mathsf K(N)\xi
 =\max_{h\in Q_NH_\nu}
       \left\{2\RePart\int_DF_\xi\overline h\dd k
                               -\mathfrak a_N(h,h)\right\}.
 \end{equation}
 The unique maximizer is $L_{N,Q}^{-1}F_\xi$.
\end{proposition}
\begin{proof}
 Equation \eqref{ons:inverse-pairing} gives
 $\mathsf K_{ij}^{ab}=\mathfrak a_N(u_i^a,u_j^b)$, and hence reciprocity.
 With $F_\xi=\sum_{j,a}\mathcal G_j^a\xi_j^a$ and
 $u_\xi=L_{N,Q}^{-1}F_\xi$,
 \begin{equation}\label{ons:quadratic-form}
 \sum_{i,j,a,b}\xi_i^a\mathsf K_{ij}^{ab}\xi_j^b
       =\mathfrak a_N(u_\xi,u_\xi)\ge0.
 \end{equation}
 Pairings of parameter derivatives are controlled by
 \eqref{ons:cell-bound} and $\nu_*^{-1}\in L^1$.
 Completing the square gives
 \[
 2\RePart\int_DF_\xi\overline h\dd k-\mathfrak a_N(h,h)
 =\mathfrak a_N(u_\xi,u_\xi)
                -\mathfrak a_N(h-u_\xi,h-u_\xi).
 \]
 The last term is nonnegative and vanishes on $Q_NH_\nu$ exactly when
 $h=u_\xi$.
\end{proof}

To relate this tensor to Euler transport, let $N=N_{\theta(t,X)}$
be a smooth positive Euler solution and set $p=1/N=\theta\cdot\Psi$.
The five-moment equations give
\begin{equation}\label{ons:euler-force}
 F_E=-\mathcal T\log N=N\mathcal Tp,
 \qquad P_NF_E=0,
 \qquad F_E=\sum_{j,a}\mathcal G_j^a\partial_j\theta_a.
\end{equation}
Indeed, $N\Psi\cdot\theta_t$ belongs to the null space, while
$\int\kappa_aF_E=-\int\Psi_a\mathcal TN=0$.
Thus, for $z=L_{N,Q}^{-1}F_E$,
\begin{equation}\label{ons:cell-flux}
 \int_Dv_i\Psi_aNz\,\dd k
 =\langle\mathcal G_i^a,z\rangle
 =\sum_{j,b}\mathsf K_{ij}^{ab}(N)\partial_j\theta_b.
\end{equation}
This identifies the coefficients and signs of all five first-order fluxes.

\subsection{Null directions in the gradient space}

Write a gradient array as $\xi=(g,B,h)$, with
$g,h\in\R^3$ and $B\in\R^{3\times3}$, according to
\[
 g_j=\xi_j^0,\qquad B_{jl}=\xi_j^l\ (1\le l\le3),
 \qquad h_j=\xi_j^4.
\]
When $\xi=\nabla_X\theta$, these are $g=\nabla_X\alpha$,
$B_{jl}=\partial_j\beta_l$ and $h=\nabla_X\gamma$.
Set
\[
 B^\circ=\frac{B+B^T}{2}-\frac{\tr B}{3}I.
\]

\begin{theorem}[Kernel and rank of the transport tensor]\label{ons:kernel}
 For every positive parameter $\theta$,
 \begin{equation}\label{ons:kernel-formula}
 \Ker\mathsf K(N)
 =\{(g,B,0):g\in\R^3,\ (B+B^T)/2=bI\text{ for some }b\in\R\}.
 \end{equation}
 The kernel has dimension $7$ and the tensor has rank $8$.
 Uniformly for $\theta\in\Theta$,
 \begin{equation}\label{ons:deviatoric-bound}
 c_\Theta(|B^\circ|^2+|h|^2)
 \le\sum_{i,j,a,b}\xi_i^a\mathsf K_{ij}^{ab}(N)\xi_j^b
 \le C_\Theta(|B^\circ|^2+|h|^2).
 \end{equation}
\end{theorem}
\begin{proof}
 By \eqref{ons:quadratic-form} and weighted coercivity, the quadratic
 form vanishes exactly when $u_\xi$ belongs to the null space.
 The five normalization conditions then give $u_\xi=0$, equivalently
 $F_\xi=0$. Dividing by the positive function $N$, this condition becomes
 \begin{equation}\label{ons:null-polynomial}
 2g\cdot k+2k^TBk+2(h\cdot k)|k|^2
       \in\Span\{1,k^{(1)},k^{(2)},k^{(3)},|k|^2\}.
 \end{equation}
 Since the cube has nonempty interior, an almost-everywhere polynomial
 identity is an identity of coefficients. The cubic terms give $h=0$,
 and the quadratic terms give $B^\circ=0$. The skew-symmetric part
 does not enter $k^TBk$, and the linear terms are already collision
 invariants. This proves \eqref{ons:kernel-formula} and gives
 dimension $3+3+1=7$.

 For the uniform comparison, $F_\xi$ depends only on the eight effective
 components $(B^\circ,h)$. On the compact set
 $|B^\circ|^2+|h|^2=1$, $\theta\in\Theta$, the five-dimensional projection
 is continuous in the parameters, and \eqref{ons:null-polynomial}
 excludes zero. Hence $\|F_\xi\|_2\ge c_\Theta>0$.
 Testing the cell equation with the bounded function $F_\xi$ and using
 Cauchy's inequality for the form yields
 \[
 \|F_\xi\|_2^4
 \le\mathfrak a_N(u_\xi,u_\xi)\mathfrak a_N(F_\xi,F_\xi)
 \le C_\Theta\mathfrak a_N(u_\xi,u_\xi)\|F_\xi\|_2^2.
 \]
 This proves the lower bound. Microscopic coercivity and weighted duality
 give the upper bound:
 \[
 \mathfrak a_N(u_\xi,u_\xi)
 \le C_\Theta\|F_\xi/\sqrt{\nu_*}\|_2^2
 \le C_\Theta(|B^\circ|^2+|h|^2).
 \]
 All integrals use the cubic domain and its four-leg resonance measure.
\end{proof}

Some pointwise gradient arrays make the null directions explicit.
Take $g=h=0$. If $B=\lambda I$, then
$2k^TBk=2\lambda|k|^2$ is a collision invariant; if $B^T=-B$, then
$k^TBk=0$. Both give zero dissipation.
In contrast, $B=\operatorname{diag}(1,-1,0)$ has a nonzero traceless
symmetric part, so \eqref{ons:deviatoric-bound} gives positive dissipation.
These are gradient arrays at a single spatial point, not global affine
parameter fields on the torus.

The degeneracy concerns gradients of the dual parameter $\theta$.
The mass component also satisfies an identity for every individual state.
If $f\in L^1(D)$ and $N$ have the same five moments, then
\begin{equation}\label{ons:mass-flux}
 \int_Dv_i(f-N)\,\dd k=2\int_Dk^{(i)}(f-N)\,\dd k=0,
 \qquad\mathcal G_i^0=0,
 \qquad\mathsf K_{ij}^{0b}=\mathsf K_{ij}^{a0}=0.
\end{equation}
The mass equation remains part of the five-moment system; momentum
matching makes its residual flux vanish.

\subsection{The spatial Fourier symbol and transport coupling}

For $\ell\in\R^3$, define
\begin{equation}\label{ons:diffusion-symbol}
 \mathsf D_N(\ell)_{ab}
       =\sum_{i,j}\ell_i\mathsf K_{ij}^{ab}(N)\ell_j.
\end{equation}
This matrix is real symmetric and nonnegative. A Fourier amplitude
$a=(a_\alpha,a_\beta,a_\gamma)$ gives the array
$\xi_j^b=\ell_ja_b$; the common factor $i$ from spatial differentiation
cancels in the quadratic form. Thus the fifteen gradient components of
a single mode cannot be chosen independently: they must have the form
$\ell\otimes a$. Substituting this constraint into
\eqref{ons:kernel-formula} reduces the seven null directions to the
single mass direction.

\begin{proposition}[The unique null direction at nonzero frequency]\label{ons:rank-one}
 If $\ell\ne0$, then
 \begin{equation}\label{ons:rank-four}
 \Ker\mathsf D_N(\ell)=\Span\{e_{\mathrm{mass}}\},\qquad
 \operatorname{rank}\mathsf D_N(\ell)=4,
 \end{equation}
 where $e_{\mathrm{mass}}=(1,0,0,0,0)^T$.
 The uniform comparison
 \begin{equation}\label{ons:directional-bound}
 c_\Theta|\ell|^2(|a_\beta|^2+|a_\gamma|^2)
 \le a^*\mathsf D_N(\ell)a
 \le C_\Theta|\ell|^2(|a_\beta|^2+|a_\gamma|^2)
 \end{equation}
 holds also for complex amplitudes.
\end{proposition}
\begin{proof}
 Theorem \ref{ons:kernel} gives $a_\gamma=0$ and
 $\operatorname{sym}(\ell\otimes a_\beta)=bI$.
 Contracting with a nonzero $z\perp\ell$ on both sides gives $b=0$.
 Contracting with $\ell$ gives $\ell\cdot a_\beta=0$, and applying
 the matrix to $\ell$ then yields $a_\beta=0$.
 The component $a_\alpha$ is arbitrary, proving the kernel statement.
 For $|\ell|=1$, take the minimum on the unit sphere of the remaining
 four-dimensional amplitude space. Equation \eqref{ons:deviatoric-bound}
 and homogeneity give the uniform comparison.
 The complex case follows by separating real and imaginary parts.
\end{proof}

Use $M,J_j$ from \eqref{hyd:matrices}. For a unit direction $\omega$, write
$J(\omega)=\sum_j\omega_jJ_j$ and
$e_{\beta\cdot\omega}=(0,\omega_1,\omega_2,\omega_3,0)^T$.
The Gram definitions give
\begin{equation}\label{ons:euler-coupling}
 J(\omega)e_{\mathrm{mass}}=2M e_{\beta\cdot\omega}.
\end{equation}
Thus no eigenvector of $M^{-1}J(\omega)$ lies in the null direction
of \eqref{ons:rank-four}. This is the Shizuta--Kawashima coupling
condition for the five-moment system \cite[Theorem 1.1]{SK1985}.
We give a direct decay proof.

Our Fourier convention is
\[
 \widehat\vartheta(\ell)=(2\pi)^{-3}
       \int_{\T^3}\vartheta(X)e^{-i\ell\cdot X}\,\dd X,
       \qquad \ell\in\mathbb Z^3.
\]
For a fixed positive matrix $M$, define
\[
 \|\vartheta\|_{H^s,M}^2
   =(2\pi)^3\sum_{\ell\in\mathbb Z^3}
      (1+|\ell|^2)^s\widehat\vartheta(\ell)^*M\widehat\vartheta(\ell).
\]
For $s=0$, Parseval's identity identifies this with
$\int_{\T^3}\vartheta^*M\vartheta\,\dd X$.
For general $s$ it is equivalent to the usual $H^s$ norm, uniformly
on the compact parameter set under consideration.

\begin{proposition}[Fourier decay of the constant-parameter system]\label{ons:fourier}
 Fix $\theta\in\Theta$ and consider on $\T^3$ the constant-coefficient system
 \begin{equation}\label{ons:frozen-system}
 M\partial_t\vartheta+\sum_iJ_i\partial_i\vartheta
       =c^{-1}\sum_{i,j}\mathsf K_{ij}\partial_i\partial_j\vartheta,
       \qquad c>0.
 \end{equation}
 For every $\ell\in\mathbb Z^3\setminus\{0\}$,
 \begin{equation}\label{ons:fourier-decay}
 \|\widehat\vartheta(t,\ell)\|_M
 \le C_\Theta\exp\!\left(-a_\Theta
           \frac{c|\ell|^2}{c^2+|\ell|^2}t\right)
           \|\widehat\vartheta(0,\ell)\|_M,
 \qquad\|z\|_M^2=z^*Mz.
 \end{equation}
 All five zero-frequency components are constant.
 For $c\ge1$ and mean-zero data, the system also satisfies, on every
 $H^s(\T^3;\R^5)$,
 \begin{equation}\label{ons:torus-decay}
 \|\vartheta(t)\|_{H^s,M}
       \le C_\Theta e^{-a_\Theta t/c}\|\vartheta(0)\|_{H^s,M}.
 \end{equation}
\end{proposition}
\begin{proof}
 Set $r=|\ell|>0$, $\omega=\ell/r$, and
 \[
 A=M^{-1/2}J(\omega)M^{-1/2},\quad
 B=M^{-1/2}\mathsf D_N(\omega)M^{-1/2},\quad
 e=\frac{M^{1/2}e_{\mathrm{mass}}}{\sqrt{M_{\alpha\alpha}}},\quad
 b=(I-ee^T)Ae.
 \]
 Both matrices are real symmetric, $Be=0$, and
 $B\ge a_0(I-ee^T)$, $\|A\|+\|B\|\le C_0$.
 By \eqref{ons:euler-coupling},
 \[
 |b|^2=\frac4{M_{\alpha\alpha}}
 \left(e_{\beta\cdot\omega}^TMe_{\beta\cdot\omega}
       -\frac{|e_{\mathrm{mass}}^TMe_{\beta\cdot\omega}|^2}
                    {M_{\alpha\alpha}}\right)\ge a_0>0.
 \]
 The expression in parentheses is a Schur complement of a positive
 Gram matrix, so the lower bound is uniform on
 $\Theta\times\mathbb S^2$.

 Let $S=be^T-eb^T$. This is a real skew-symmetric matrix and
 $e^T[S,A]e=-2|b|^2$. For a fixed sufficiently small $\eta>0$, set
 \[
 d=r/c,\qquad \alpha_d=\frac{\eta d}{1+d^2},\qquad
 H=I+i\alpha_dS.
 \]
 Then $H$ is Hermitian and $I/2\le H\le3I/2$.
 For $y=M^{1/2}\widehat\vartheta$, the equation reads
 $y'=\mathcal Gy$, with $\mathcal G=-irA-c^{-1}r^2B$.
 Direct multiplication gives
 \begin{equation}\label{ons:compensated-matrix}
 H\mathcal G+\mathcal G^*H
 =-2c^{-1}r^2B+\alpha_dr[S,A]
       -i\alpha_dc^{-1}r^2(SB+BS).
 \end{equation}
 Write $y=ze+q$, $q\perp e$. Since $Be=0$, the matrix $SB+BS$
 has only cross blocks between $e$ and $e^\perp$.
 Hence $E=y^*Hy$ satisfies
 \[
 E'\le-a_1c^{-1}r^2|q|^2-a_1\alpha_dr|z|^2
       +C\alpha_dr|q|^2+C\alpha_dr(1+d)|z||q|.
 \]
 Young's inequality absorbs the mixed term using half of
 $\alpha_dr|z|^2$, at the cost of at most
 \[
 C\alpha_dr(1+d^2)|q|^2=C\eta c^{-1}r^2|q|^2.
 \]
 Choosing $\eta$ sufficiently small yields
 \[
 E'\le-a_2\{c^{-1}r^2|q|^2+\alpha_dr|z|^2\}
       \le-a_3\frac{cr^2}{c^2+r^2}E.
 \]
 Integration and the comparison of $H$ with $I$ give
 \eqref{ons:fourier-decay}. At $\ell=0$, the matrix equation is
 $M\partial_t\widehat\vartheta=0$.
 If $c\ge1$ and $|\ell|\ge1$, then
 $c|\ell|^2/(c^2+|\ell|^2)\ge1/(2c)$.
 Summing over Fourier modes gives \eqref{ons:torus-decay}.
 Density of finite Fourier sums and continuity of each mode also
 give a strongly continuous semigroup on $H^s$.
\end{proof}

The tensor in this section uses the $c$ normalization in
\eqref{hyd:kinetic}. With $\varepsilon=4\pi/c$, the tensor multiplying
the first-order parameter $\varepsilon$ is $\mathsf K/(4\pi)$:
\[
 c^{-1}\mathsf K=\varepsilon\,\frac{\mathsf K}{4\pi}.
\]
